
\subsubsection{2.1 Spurious Correlation Robustness}

Sagawa et al. [2019] formalize spurious correlation robustness as a benchmark challenge with Waterbirds and the GroupDRO objective. GroupDRO achieves strong WGA by minimizing the maximum group-specific loss at each step but requires group annotations for all training samples. Idrissi et al. [2021] demonstrate that simple data balancing over groups is a competitive baseline on these benchmarks, provided group identity is known, underscoring that the core challenge is identifying which samples belong to minority groups.

\subsubsection{2.2 Label-Free Two-Stage Methods}

\textbf{Just Train Twice (JTT)} \citep{Liu2021Just} identifies ERM model misclassifications as a minority proxy and upweights this error set in a second training stage. The proxy signal is binary and coarse: a sample is either misclassified or not.

\textbf{Learning from Failure (LfF)} \citep{Nam2020Learning} trains two networks simultaneously and uses the ratio of generalized cross-entropy losses to identify hard samples for upweighting. It requires two parallel training streams.

\textbf{Deep Feature Reweighting (DFR)} \citep{Kirichenko2022Last} demonstrates that ERM-trained features already encode class-relevant information and that the problem resides in the classifier head. DFR freezes the ERM backbone and retrains only the last layer using a group-balanced validation subset. This achieves strong WGA on Waterbirds but requires group-labeled samples for subset construction. The proposed Stage 2 of the present pipeline follows the DFR principle while replacing group-labeled subset construction with gradient-norm-based selection. Stage 2 was not executed in the reported experiments.

No prior work in this paradigm uses per-sample gradient norms as the minority identification signal.

\subsubsection{2.3 Gradient Dynamics and Related Signals}

The \textbf{EL2N score} \citep{Paul2021} is defined as the expected L2 norm of the prediction-error vector, $E\|p_i - y_{i}\|$. The normalized gradient norm $\tilde{g}_i = \|p_i - y_{i}\|$ is mathematically equivalent to EL2N evaluated at a single checkpoint. EL2N was introduced for dataset pruning by selecting easy (low-error) samples; the present work applies the same per-sample signal in the opposite selection direction --- to identify hard (high-error) minority samples --- for a different purpose.

\textbf{Forgetting events} \citep{Toneva2019} and \textbf{influence functions} \citep{Koh2017} are training-dynamic characterizations of individual samples but have not been applied to minority group identification in spuriously correlated settings.

\textbf{GEORGE} \citep{Zhang2022GEORGE} discovers pseudo-groups by unsupervised clustering in the model's representation space, then applies GroupDRO. This requires a full clustering step over the training set. The present method uses a scalar per-sample signal requiring no clustering.

\textbf{Edge of Stability} dynamics \citep{Cohen2021Gradient} establish that gradient noise near the edge-of-stability regime can amplify per-sample learning rate differences, which may contribute to gradient norm disparity between groups.

\subsubsection{2.4 Position}

The normalized gradient norm $\tilde{g}$ provides a continuous minority proxy signal derived entirely from last-layer training dynamics, requiring no change to the training objective, no parallel networks, no clustering, and no group labels. The outer-product decomposition of the cross-entropy gradient provides a mechanistic interpretation that is absent from other proxy signals.

